\begin{abstract}
Large language models hallucinate---generating confident but incorrect statements---yet current detection methods require multiple forward passes (semantic entropy) or external knowledge bases. We propose Cross-Layer Trajectory Instability (CLTI), a single-pass approach measuring how representations evolve across late transformer layers. On TruthfulQA MC1 with LLaMA-2-7B (layers 24--31), Normalized Trajectory Instability (NTI) achieves AUROC 0.5657; combined with Convergence Monotonicity Index (CMI) and output entropy, detection improves by $+7.1\%$ AUROC ($p = 1.15 \times 10^{-5}$) over an intercept-only baseline. However, trajectory metrics fail on low-entropy (``confident'') predictions (AUROC 0.5136), and Representational Competition Index (RCI) flip patterns appear in $>90\%$ of all responses, revealing architectural rather than epistemic dynamics. Our findings demonstrate that trajectory features provide complementary signal for uncertain predictions, while identifying limitations that constrain the method's scope. We release code for trajectory feature extraction from any transformer.
\end{abstract}
